
%01. The Reality Beyond the Illusion of Space-Time.tex

\section{애초에 시간은 아니었다 - 우주의 진짜 하드웨어 평활 엔진}
우리는 시간이 흐른다고 믿는다. 우리는 삶의 모든 사건들이 시간이라는 보이지 않는 실 위에 나란히 배열되어 있다고 생각한다. 마치 목걸이에 보석들이 하나씩 꿰어져 있는 것처럼 말이다. 과거는 우리 뒤에 놓여 있고 미래는 우리 앞에 펼쳐져 있으며, 각각의 사건은 그 실 위에서 저마다의 위치를 차지한 채 서로 측정 가능한 간격을 두고 존재하는 것처럼 보인다. 아침에 눈을 뜨고, 시계를 확인하고, 약속 시간에 맞추어 이동하며, 스마트폰 화면에는 초 단위로 시간이 표시된다. 하지만 어느 순간부터 나는 이런 생각을 하게 되었다. 정말 시간이 먼저 존재하고, 그 시간 속에서 사건이 발생하는 것일까? 공간이 먼저 존재하고 그 속에서 물체가 움직이는 것일까? 아니면 우리가 시간과 공간이라고 부르는 것은 어떤 더 근본적인 작용을 설명하기 위해 인간이 나중에 붙인 이름일 뿐일까.

나는 시간, 공간, 거리, 속도가 우주의 하드웨어가 아니라, 인간이 현실을 다루기 위해 만들어낸 '인터페이스'에 가깝다고 생각한다. 컴퓨터 화면 속 폴더와 휴지통이 내부의 복잡한 전기적 신호를 인간이 다루기 쉽게 만든 상징인 것처럼, 시간과 공간 역시 우리의 생존과 인식을 위해 구성된 화면일 수 있다. 우리의 뇌는 우주의 근본 구조를 이해하기 위해서가 아니라, 던져진 돌의 궤적을 예측하고 맹수와의 거리를 판단하기 위해 진화했기 때문이다.

그렇다면 이 시간과 공간이라는 화면 뒤에 있는 실제 하드웨어는 무엇인가. 나는 그 근본에 '에너지 갭(Energy Gap)'이 있다고 생각한다. 산꼭대기에 있는 바위가 굴러떨어지고 계곡물이 흐르는 사건의 근본 원인은 시간이나 공간이 아니라, 높은 곳과 낮은 곳 사이의 포텐셜 에너지 차이, 즉 에너지 갭이 존재하고 그것이 해소되고 있기 때문이다. 이 관점에서 보면 빅뱅 역시 빈 공간 속에서 물질이 폭발한 사건이라기보다, 상상할 수 없는 수준의 에너지 집중 상태가 만들어낸 최대의 에너지 갭으로 이해될 수 있다. 우주는 거의 영원이라는 시간 동안 단순히 팽창한 것이 아니라, 최초의 극단적인 에너지 갭을 해소하는 방향으로 끊임없이 '떨어지고' 있는 것이다.

여기서 나는 엔트로피의 의미도 다시 생각하게 되었다. 흔히 엔트로피를 무질서나 혼돈이라고 부르지만, 나는 이를 '평탄화(Equalization)'로 보고 싶다. 유용한 에너지 갭이 사라져 흐름이 멈추고, 더 이상 어떤 차이도 남아 있지 않은 완전한 평탄 상태에 가까워지는 것이 바로 엔트로피가 최대가 되는 상태이다.

그렇다면 우주의 근본 경향이 모든 것이 평탄해지는 방향으로 진행되는 것인데, 왜 생명처럼 복잡하고 역동적인 구조가 생겨났을까? 오랫동안 생명은 외부에서 음의 엔트로피를 끌어와 자기 내부의 질서를 유지하며 우주적 흐름에 저항하는 예외적인 존재로 이해되어 왔다. 하지만 나는 여기서 한 걸음 더 나아가야 한다고 생각한다. 생명이 유지하는 국소적 질서는 우주에 저항하기 위해서가 아니라, 오히려 더 큰 엔트로피 증가를 위한 장치적 조건일 뿐이다.

우리가 밥을 먹고 소화하며 몸을 움직이는 과정은 고밀도 에너지를 끌어와 저밀도 에너지로 바꾸고 주변에 흩뿌리는 행위이다. 겉으로 보면 생명체는 자기 몸의 질서를 유지하는 것 같지만, 실상은 그 작은 질서를 유지하기 위해 주변 세계에 훨씬 더 큰 에너지 분산을 일으킨다. 이렇게 보면 생명은 우주에 저항하는 고귀한 예외가 아니라, 우주의 평탄화를 맹렬히 가속하는 '로컬 엔진(Local Engine)'이다. 먹고, 소화하고, 움직이고, 번식하는 모든 과정은 에너지 갭을 해소하는 과정의 일부인 것이다. 나는 생명이 우주의 법칙 바깥에 있는 것이 아니라, 그 법칙이 만들어낸 고도의 에너지 분산 구조라고 본다.

